% GENERATED FILE. Edit canonical lessons, then run npm run generate:books.
% canonical-source-hash: fnv1a64:9c55c00334fb17bb
% canonical-lessons: KA-C128-mimcu, KA-C128-birugali, KA-C128-ale, KA-C128-jalapata, KA-C128-dvipa

\chapter{Lightning, Storm, Wave, Waterfall, Island}
\label{ch:ka-lightning-storm-wave-waterfall-island}
\hlchaptermodality{128}

\begin{chapteropening}
\textbf{By the end of this chapter:} \emph{I can say lightning, a storm, a wave, a waterfall and an island.}

Complete the last lesson of chapter 128: I can say lightning, a storm, a wave, a waterfall and an island.
\end{chapteropening}

\section[\texorpdfstring{\kn{ಮಿಂಚು} (miṁcu)}{ಮಿಂಚು (miṁcu)}]{\kn{ಮಿಂಚು} (miṁcu) — lightning}

\label{lesson:KA-C128-mimcu}



\begin{quote}
Before the new one: say the Kannada for a stone, then the Kannada for thunder.
\end{quote}

\subsection*{You'll want to know: \kn{ಮಿಂಚು}}
\textbf{\kn{ಮಿಂಚು}} — \emph{miṁcu} — \textquotedblleft{}lightning\textquotedblright{}.

\begin{grammarlens}[title={What you've built}]
One more word for this chapter.
\end{grammarlens}

\subsection*{Your turn}
\begin{itemize}
\raggedright
  \item \emph{Say it:} \emph{miṁcu}
  \item \emph{Say it:} \emph{miṁcu}, once more
  \item \emph{Say it:} say \emph{miṁcu}
  \item \emph{Recall:} say the Kannada for a stone, then the Kannada for thunder, then say \emph{miṁcu} again
\end{itemize}

\begin{tcolorbox}[breakable,colback=teal!4,colframe=teal!35!black,title={Before you move on}]
What does \kn{ಮಿಂಚು} mean? (\textquotedblleft{}Lightning\textquotedblright{}.) Say it once more.
\end{tcolorbox}
\section[\texorpdfstring{\kn{ಬಿರುಗಾಳಿ} (birugāḷi)}{ಬಿರುಗಾಳಿ (birugāḷi)}]{\kn{ಬಿರುಗಾಳಿ} (birugāḷi) — a storm}

\label{lesson:KA-C128-birugali}



\begin{quote}
Before the new one: say the Kannada for thunder, then the Kannada for lightning.
\end{quote}

\subsection*{You'll want to know: \kn{ಬಿರುಗಾಳಿ}}
\textbf{\kn{ಬಿರುಗಾಳಿ}} — \emph{birugāḷi} — \textquotedblleft{}a storm\textquotedblright{}.

\begin{grammarlens}[title={What you've built}]
One more word for this chapter.
\end{grammarlens}

\subsection*{Your turn}
\begin{itemize}
\raggedright
  \item \emph{Say it:} \emph{birugāḷi}
  \item \emph{Say it:} \emph{birugāḷi}, once more
  \item \emph{Say it:} say \emph{birugāḷi}
  \item \emph{Recall:} say the Kannada for thunder, then the Kannada for lightning, then say \emph{birugāḷi} again
\end{itemize}

\begin{tcolorbox}[breakable,colback=teal!4,colframe=teal!35!black,title={Before you move on}]
What does \kn{ಬಿರುಗಾಳಿ} mean? (\textquotedblleft{}A storm\textquotedblright{}.) Say it once more.
\end{tcolorbox}
\section[\texorpdfstring{\kn{ಅಲೆ} (ale)}{ಅಲೆ (ale)}]{\kn{ಅಲೆ} (ale) — a wave}

\label{lesson:KA-C128-ale}



\begin{quote}
Before the new one: say the Kannada for lightning, then the Kannada for a storm.
\end{quote}

\subsection*{You'll want to know: \kn{ಅಲೆ}}
\textbf{\kn{ಅಲೆ}} — \emph{ale} — \textquotedblleft{}a wave\textquotedblright{}.

\begin{grammarlens}[title={What you've built}]
One more word for this chapter.
\end{grammarlens}

\subsection*{Your turn}
\begin{itemize}
\raggedright
  \item \emph{Say it:} \emph{ale}
  \item \emph{Say it:} \emph{ale}, once more
  \item \emph{Say it:} say \emph{ale}
  \item \emph{Recall:} say the Kannada for lightning, then the Kannada for a storm, then say \emph{ale} again
\end{itemize}

\begin{tcolorbox}[breakable,colback=teal!4,colframe=teal!35!black,title={Before you move on}]
What does \kn{ಅಲೆ} mean? (\textquotedblleft{}A wave\textquotedblright{}.) Say it once more.
\end{tcolorbox}
\section[\texorpdfstring{\kn{ಜಲಪಾತ} (jalapāta)}{ಜಲಪಾತ (jalapāta)}]{\kn{ಜಲಪಾತ} (jalapāta) — a waterfall}

\label{lesson:KA-C128-jalapata}



\begin{quote}
Before the new one: say the Kannada for a storm, then the Kannada for a wave.
\end{quote}

\subsection*{You'll want to know: \kn{ಜಲಪಾತ}}
\textbf{\kn{ಜಲಪಾತ}} — \emph{jalapāta} — \textquotedblleft{}a waterfall\textquotedblright{}.

\begin{grammarlens}[title={What you've built}]
One more word for this chapter.
\end{grammarlens}

\subsection*{Your turn}
\begin{itemize}
\raggedright
  \item \emph{Say it:} \emph{jalapāta}
  \item \emph{Say it:} \emph{jalapāta}, once more
  \item \emph{Say it:} say \emph{jalapāta}
  \item \emph{Recall:} say the Kannada for a storm, then the Kannada for a wave, then say \emph{jalapāta} again
\end{itemize}

\begin{tcolorbox}[breakable,colback=teal!4,colframe=teal!35!black,title={Before you move on}]
What does \kn{ಜಲಪಾತ} mean? (\textquotedblleft{}A waterfall\textquotedblright{}.) Say it once more.
\end{tcolorbox}
\section[\texorpdfstring{\kn{ದ್ವೀಪ} (dvīpa)}{ದ್ವೀಪ (dvīpa)}]{\kn{ದ್ವೀಪ} (dvīpa) — an island}

\label{lesson:KA-C128-dvipa}



\begin{quote}
Before the new one: say the Kannada for a wave, then the Kannada for a waterfall.
\end{quote}

\subsection*{You'll want to know: \kn{ದ್ವೀಪ}}
\textbf{\kn{ದ್ವೀಪ}} — \emph{dvīpa} — \textquotedblleft{}an island\textquotedblright{}.

\begin{grammarlens}[title={What you've built}]
One more word for this chapter.
\end{grammarlens}

\subsection*{Your turn}
\begin{itemize}
\raggedright
  \item \emph{Say it:} \emph{dvīpa}
  \item \emph{Say it:} \emph{dvīpa}, once more
  \item \emph{Say it:} say \emph{dvīpa}
  \item \emph{Recall:} say the Kannada for a wave, then the Kannada for a waterfall, then say \emph{dvīpa} again
\end{itemize}

\begin{tcolorbox}[breakable,colback=teal!4,colframe=teal!35!black,title={Before you move on}]
What does \kn{ದ್ವೀಪ} mean? (\textquotedblleft{}An island\textquotedblright{}.) Say it once more.
\end{tcolorbox}
